\documentclass[10pt]{article}
\usepackage[margin=1in]{geometry}
\usepackage{setspace}
\singlespacing
\usepackage{enumitem}
\usepackage{booktabs}
\usepackage{graphicx}
\usepackage{hyperref}
\usepackage{xcolor}
\usepackage{listings}
\usepackage{tabularx}

\lstset{
  basicstyle=\ttfamily\scriptsize,
  breaklines=true,
  frame=single,
  backgroundcolor=\color{gray!10},
}

\setlength{\parskip}{4pt}
\setlength{\parindent}{0pt}

\title{\textbf{Spec-to-Tapeout: An LLM Multi-Agent Pipeline\\for Automated ASIC Design}}
\author{Charlie Wu \\ EEE 598 --- ASU ICLAD 2025 Hackathon}
\date{}

\begin{document}
\maketitle

\section{Methodology \& Problem Definition}

The ICLAD Spec-to-Tapeout challenge requires building an automated pipeline that reads a YAML hardware specification and produces a tapeout-ready ASIC design targeting the SkyWater 130nm HD (sky130hd) process node.

\textbf{Inputs:} YAML files specifying a module name, port definitions (direction, type, width), design type (FSM, pipelined, or combinational), clock period, and optional parameters.

\textbf{Outputs:} Per problem: (1) synthesizable SystemVerilog RTL that passes functional verification via Icarus Verilog, and (2) a final OpenROAD database (\texttt{6\_final.odb}) produced by running OpenROAD-flow-scripts (ORFS) through synthesis, placement, clock-tree synthesis, routing, and finishing.

\textbf{Objective:} Maximize a 100-point score computed from four physical metrics compared against a reference design: worst negative slack (WNS, 40\,pts), total negative slack (TNS, 20\,pts), total power (20\,pts), and instance area (20\,pts).

\textbf{Approach:} Rather than a single monolithic LLM prompt, the system uses a multi-agent pipeline with specialized stages---specification parsing, multi-strategy RTL generation, iterative verification and repair, physical synthesis, metric-based ranking, and targeted optimization. This decomposition lets each stage use focused prompts with domain-specific constraints, dramatically improving success rates over one-shot generation.

\section{Agent Architecture}

The pipeline consists of seven cooperating agents orchestrated by an async Python runtime (\texttt{spec2tapeout\_agent.py}):

\begin{enumerate}[leftmargin=*,topsep=2pt,itemsep=1pt]
  \item \textbf{Spec Interpreter} --- Parses YAML into a structured \texttt{Spec} object. Handles edge cases such as Python-like syntax in YAML fields (e.g., \texttt{[0,1]*40} alias conflicts) by stripping non-essential sections.
  \item \textbf{RTL Generator} --- Produces three independent RTL candidates per problem using distinct strategies: \emph{textbook} (clean, correct), \emph{timing\_opt} (minimize critical path), and \emph{area\_opt} (minimize gate count). Each strategy injects different design guidance into the LLM prompt.
  \item \textbf{Verifier} --- Compiles each candidate with \texttt{iverilog -Wall -g2012} against the problem testbench and runs simulation with \texttt{vvp}. Checks for PASSED/FAILED in the output.
  \item \textbf{RTL Fixer} --- When verification fails, sends the error log, original spec, testbench source, and current RTL back to the LLM for targeted repair. Supports up to 3 fix-and-retry iterations per candidate, with separate handling for compilation vs.\ simulation errors.
  \item \textbf{ORFS Runner} --- Executes the full OpenROAD flow inside a Docker container (\texttt{openroad/flow-ubuntu22.04-builder:836842}). Generates \texttt{config.mk} and \texttt{constraint.sdc} from the spec, mounts the workspace, and runs \texttt{make finish}. Includes automatic retry with relaxed utilization on failure.
  \item \textbf{Ranker} --- Reads \texttt{6\_report.json} produced by ORFS, strips the \texttt{finish\_\_} key prefix, and scores each candidate against reference metrics using the official scoring formula.
  \item \textbf{Optimizer} --- If the best score falls below 85/100, the LLM receives the current RTL along with specific metrics feedback (e.g., ``TIMING CRITICAL: WNS=$-$1.47ns'') and attempts a targeted optimization pass.
\end{enumerate}

\textbf{Difference from simple prompting:} A one-shot prompt produces a single candidate with no feedback. This system generates multiple diverse candidates, verifies each against ground-truth testbenches, iteratively repairs failures, evaluates physical metrics after real EDA synthesis, and closes the loop with optimization. The multi-strategy approach ensures coverage---for example, P1 (FSM) scores highest with \emph{textbook} while P7 (Taylor series) benefits from \emph{optimized}.

All five problems execute in parallel via \texttt{asyncio.gather}, with ORFS concurrency limited to 3 simultaneous Docker containers.

\section{Tools \& Frameworks}

\begin{table}[h]
\centering
\small
\begin{tabularx}{\textwidth}{lX}
\toprule
\textbf{Tool} & \textbf{Role} \\
\midrule
OpenAI API (GPT-5.3-codex, GPT-5.4) & RTL generation, error repair, optimization \\
Icarus Verilog (\texttt{iverilog}/\texttt{vvp}) & Functional verification (compile + simulate) \\
OpenROAD-flow-scripts (ORFS) via Docker & Full physical synthesis (Yosys $\to$ floorplan $\to$ place $\to$ CTS $\to$ route $\to$ finish) \\
Yosys (inside ORFS) & Logic synthesis (RTL $\to$ gate-level netlist) \\
Python 3 + asyncio & Pipeline orchestration, async parallel execution \\
PyYAML & YAML specification parsing \\
Docker & Containerized EDA environment \\
\bottomrule
\end{tabularx}
\end{table}

No external agent framework is used; the orchestration is implemented directly in Python with \texttt{asyncio} for parallelism.

\section{Context Engineering}

Each LLM call is constructed with carefully scoped context:

\textbf{System prompt} includes: (1) base rules (synthesizable-only constructs, exact module signature matching), (2) Yosys compatibility constraints (no blocking assigns in \texttt{always\_ff}, default assignments in \texttt{always\_comb} to prevent latch inference, no nested always blocks), (3) strategy-specific guidance, and (4) design-type hints (e.g., for pipelined designs: unsigned arithmetic for packed array indexing, tree-structured accumulation).

\textbf{User prompt} includes: module name, description, clock period, port definitions, the exact module signature, and sample I/O when available.

\textbf{Fixer context} includes: the error log, original spec, testbench source (read-only), and current RTL source.

\textbf{Optimizer context} includes: current score, specific metrics (WNS, TNS, power, area), computed focus areas (e.g., ``TIMING CRITICAL''), and current RTL.

Context is managed by limiting each prompt to the minimum information needed for its task. Testbench source is included only in fixer prompts. ORFS metrics are only passed during optimization.

\section{Interaction with EDA Tools}

\textbf{Icarus Verilog:} The agent writes the candidate RTL to a temporary file, compiles it against the testbench with \texttt{iverilog -Wall -g2012 -o sim.vvp design.v testbench.v}, and runs \texttt{vvp sim.vvp}. The simulation output is searched for ``PASS'' or ``FAIL'' strings.

\textbf{ORFS (Docker):} The agent generates two EDA configuration files from the YAML spec:
\begin{itemize}[topsep=2pt,itemsep=0pt]
  \item \texttt{config.mk} --- sets \texttt{DESIGN\_NAME}, \texttt{PLATFORM=sky130hd}, \texttt{VERILOG\_FILES}, \texttt{SDC\_FILE}, utilization, and density parameters.
  \item \texttt{constraint.sdc} --- \texttt{create\_clock} with the spec's period, input/output delays at 20\% of clock period.
\end{itemize}
These are mounted into the Docker container along with the RTL, and the flow runs \texttt{make DESIGN\_CONFIG=/design/config.mk finish}. Results and logs are mounted back to the host.

\textbf{Metric extraction:} After ORFS completes, the ranker reads \texttt{6\_report.json} from the logs directory, strips the \texttt{finish\_\_} prefix from keys, and extracts \texttt{timing\_\_setup\_\_ws}, \texttt{timing\_\_setup\_\_tns}, \texttt{power\_\_total}, and \texttt{design\_\_instance\_\_area}. This approach avoids running a separate scoring container---it reads the metrics ORFS already computed.

\textbf{Iterative feedback:} The pipeline operates as a closed loop. If verification fails, the error is fed back to the LLM for repair (up to 3 attempts). If no candidate passes, an emergency regeneration step includes all prior error logs to avoid repeating mistakes. If the physical score is below threshold, the optimizer receives specific metric feedback.

\section{Prototype Implementation}

\textbf{Fully functional:}
\begin{itemize}[topsep=2pt,itemsep=0pt]
  \item End-to-end pipeline: YAML $\to$ RTL $\to$ verification $\to$ ORFS $\to$ scoring $\to$ optimization
  \item Multi-strategy generation (3 candidates per problem)
  \item Iterative verification with LLM-based repair (3 retries)
  \item Parallel execution of all 5 problems
  \item Automatic Yosys error avoidance via prompt engineering
  \item Workspace cleanup between runs (Docker-based, handles root-owned files)
  \item Generalization to hidden problems (same YAML schema)
\end{itemize}

\textbf{Limitations:}
\begin{itemize}[topsep=2pt,itemsep=0pt]
  \item No Yosys pre-synthesis check before running the full ORFS flow---Yosys failures waste an entire ORFS run (~10 min).
  \item Optimization is limited to a single pass; iterative PPA refinement would improve scores.
  \item The agent cannot directly debug timing violations at the gate level---it relies on the LLM to infer structural changes from high-level metric feedback.
\end{itemize}

\section{Initial Results}

All five visible problems produce valid solutions. Scoring is on a 100-point scale.

\begin{table}[h]
\centering
\small
\begin{tabular}{llrrrrrr}
\toprule
\textbf{Prob.} & \textbf{Module} & \textbf{Score} & \textbf{WNS (ns)} & \textbf{TNS (ns)} & \textbf{Power (W)} & \textbf{Area} \\
\midrule
P1 & seq\_detector\_0011  & \textbf{94.1} & 0.155  &  0.0  & 0.00037 & 213 \\
P5 & dot\_product         & 75.9          & 0.128  &  0.0  & 0.01625 & 35394 \\
P7 & exp\_fixed\_point    & 82.3          & 0.072  &  0.0  & 0.00163 & 9189 \\
P8 & fp16\_multiplier     & 62.9          & $-$2.012 & $-$27.7 & 0.00437 & 7983 \\
P9 & fir\_filter          & 65.0          & $-$1.888 & $-$43.2 & 0.15582 & 104787 \\
\midrule
\multicolumn{2}{l}{\textbf{Average}} & \textbf{76.0} & & & & \\
\bottomrule
\end{tabular}
\end{table}

\begin{table}[h]
\centering
\small
\begin{tabular}{lrrrr}
\toprule
\textbf{Prob.} & \textbf{Ref WNS} & \textbf{Ref TNS} & \textbf{Ref Power} & \textbf{Ref Area} \\
\midrule
P1 & 0.001 & 0.0 & 0.00060 & 244 \\
P5 & 0.230 & 0.0 & 0.03363 & 25107 \\
P7 & 0.105 & 0.0 & 0.00933 & 8539 \\
P8 & 0.260 & 0.0 & 0.00364 & 6484 \\
P9 & $-$0.166 & $-$0.7 & 0.10310 & 104542 \\
\bottomrule
\end{tabular}
\caption{Reference design metrics for comparison.}
\end{table}

P1 and P7 meet timing (WNS $>$ 0) with strong scores. P5 meets timing and achieves reasonable PPA. P8 and P9 have timing violations; P8's combinational FP16 multiplier has a critical path exceeding the 9\,ns clock, and P9's FIR filter accumulation is too deep for the 8\,ns clock.

\section{AI Usage \& Insights}

\textbf{Role of AI:} The LLM (GPT-5.3-codex) is the core generative component. It performs: (1) RTL code generation from natural-language specs, (2) targeted error repair from compiler/simulator logs, and (3) PPA-guided optimization from physical metrics.

\textbf{Where AI excels:}
\begin{itemize}[topsep=2pt,itemsep=0pt]
  \item \emph{FSM design} (P1): The LLM produces near-optimal state machines consistently, scoring 94/100 with minimal fix iterations.
  \item \emph{Error repair}: When given a compilation error or simulation mismatch along with the testbench, the LLM typically fixes the issue within 1--2 attempts.
  \item \emph{Generalization}: The same prompts and pipeline handle FSMs, pipelined datapaths, and combinational logic without problem-specific tuning.
\end{itemize}

\textbf{Where AI struggles:}
\begin{itemize}[topsep=2pt,itemsep=0pt]
  \item \emph{Timing-aware design}: The LLM has no model of gate delays at 130\,nm. It cannot predict whether a combinational path will meet a clock constraint. P8 (FP16 multiply) consistently violates timing because the LLM doesn't know that an 11$\times$11 multiply + normalize + round exceeds 9\,ns in sky130hd.
  \item \emph{Yosys compatibility}: iverilog and Yosys accept different SystemVerilog subsets. Code that passes iverilog may fail Yosys synthesis (e.g., blocking assignments in \texttt{always\_ff}, latch inference from incomplete \texttt{always\_comb}). This required explicit Yosys rules in the system prompt.
  \item \emph{Signed/unsigned semantics}: iverilog's treatment of packed array indexing (drops signedness) caused persistent failures in P5 (dot product) until the prompt explicitly warned against \texttt{\$signed()} casts.
\end{itemize}

\section{Challenges, Limitations, \& Improvements}

\textbf{Key challenges encountered:}
\begin{itemize}[topsep=2pt,itemsep=0pt]
  \item \emph{Scoring reliability}: The initial scoring approach ran a separate Docker container to extract metrics via OpenROAD, but this produced different results than the ORFS flow itself. Switching to read \texttt{6\_report.json} directly resolved inconsistencies.
  \item \emph{ORFS stale artifacts}: Docker creates files as root. When re-running the pipeline, ORFS's incremental build reused old synthesis results even with new RTL. Fixed by cleaning results/logs directories before each run using a Docker-based cleanup.
  \item \emph{LLM backend migration}: Depleted API credits required switching from Anthropic SDK to OpenAI Responses API, necessitating abstraction of the LLM client.
\end{itemize}

\textbf{Limitations:}
\begin{itemize}[topsep=2pt,itemsep=0pt]
  \item Stochastic output: LLM-generated RTL varies between runs. The same pipeline can score 86 or 70 for P7 depending on the quality of the generated code.
  \item No gate-level feedback: The agent cannot inspect the synthesized netlist to identify specific critical paths for targeted restructuring.
  \item Single optimization pass: More iterative refinement (multiple rounds of ORFS + feedback) would improve scores.
\end{itemize}

\textbf{Planned improvements for Phase 3:}
\begin{itemize}[topsep=2pt,itemsep=0pt]
  \item Add a Yosys pre-synthesis check before running the full ORFS flow to catch synthesis errors in seconds instead of minutes.
  \item Implement multi-round optimization with gate-level timing path extraction to give the LLM specific critical-path information.
  \item Increase candidate count from 3 to 5+ for more diverse coverage.
  \item Add design-specific templates (e.g., a known-good FP16 multiplier skeleton) as few-shot examples in the prompt.
\end{itemize}

\end{document}
